%% crpsp-v2-doc.tex
%% Documentação do sistema editorial LaTeX v2 do CRP-SP,
%% no estilo do pacote `doc` (classe ltxdoc).
%% Compilar 2x com: lualatex crpsp-v2-doc.tex
%% (não requer \DocumentMetadata: este manual não é PDF/UA-2)
\documentclass{ltxdoc}
\usepackage{iftex}
\ifluatex\usepackage{fontspec}\else\usepackage[T1]{fontenc}\fi
\usepackage{xcolor}
\usepackage{hyperref}
\hypersetup{colorlinks, linkcolor=blue!50!black, urlcolor=blue!50!black,
  pdftitle={Sistema editorial LaTeX v2 do CRP-SP}}

\EnableCrossrefs
\CodelineIndex
\RecordChanges

\newcommand{\pkg}[1]{\textsf{#1}}
\newcommand{\cls}[1]{\textsf{#1}}
\newcommand{\opt}[1]{\texttt{#1}}

\title{Sistema editorial \LaTeX\ v2 do CRP-SP\\
  \large\pkg{crpsp-base} \textbullet\ \cls{crpsp\_acessivel} \textbullet\
  \cls{crpsp-guia\_visual} \textbullet\ \pkg{guia-visual}}
\author{Conselho Regional de Psicologia de São Paulo\\
  \texttt{editorial/latex\_acessivel/desenvolvimento/v2/}}
\date{2026-07-03 \quad v0.1/v0.2}

\begin{document}
\maketitle

\begin{abstract}
Este manual documenta a segunda geração do sistema de classes \LaTeX\ do
CRP-SP para publicações institucionais acessíveis em conformidade
PDF/UA-2, usando o mecanismo de \emph{tagging} automático do
\pkg{latex-lab} (\texttt{testphase=\{phase-III\}}). A arquitetura separa
a infraestrutura comum (\pkg{crpsp-base}) das classes por formato
(\cls{crpsp\_acessivel} para A5, \cls{crpsp-guia\_visual} para 16:9) e
das implementações de layout por linha editorial
(\pkg{guia-crpsp\_acessivel}, \pkg{guia-visual}).
\end{abstract}

\tableofcontents

\section{Requisitos e compilação}

Todos os arquivos exigem \textbf{LuaLaTeX}; outras engines abortam com
erro. Neste workspace a compilação é feita via podman:

\begin{verbatim}
sh .../docker/lualatex-podman.sh arquivo.tex
\end{verbatim}

Compile \textbf{duas vezes} (três na primeira vez): a barra de navegação
da linha 16:9 usa um stream |.nav| dedicado (seção~\ref{sec:nav}) e o
posicionamento TikZ usa \opt{remember picture}, ambos resolvidos apenas
na passagem seguinte.

\DescribeMacro{\DocumentMetadata}
Todo documento de produção \textbf{deve} começar com o bloco abaixo,
\emph{antes} de |\documentclass| (o kernel não permite que a classe o
emita):
\begin{verbatim}
\DocumentMetadata{
  lang        = pt-BR,
  pdfversion  = 2.0,
  pdfstandard = ua-2,
  testphase   = {phase-III},
}
\end{verbatim}
Sem ele o documento compila, mas sem árvore de tags — os comandos de
tagging dos pacotes viram no-ops silenciosos (via \pkg{tagpdf-base}).

\section{\pkg{crpsp-base} — infraestrutura compartilhada}

Pacote carregado por ambas as classes. Não é usado diretamente pelo
autor, mas define a interface pública de cores.

\subsection{Verificação de engine}
Aborta a compilação com |\PackageError| se a engine não for LuaLaTeX.

\subsection{Detecção de fontspec}
\DescribeMacro{\ifcrpsp@fontspec}
Testa em Lua se \pkg{luaotfload} está funcional e define o condicional
interno |\ifcrpsp@fontspec|. Os pacotes de linha editorial usam esse
condicional para escolher entre |\setmainfont| (OpenType) e o fallback
NFSS com codificação T1. O probe aceita tanto o retorno-tabela quanto o
retorno-|true| de |require| (módulo já carregado no format, caso do
Lua\LaTeX\ moderno).

\subsection{Paleta de cores por publicação}

\DescribeMacro{\crpspCorPrincipal}
\DescribeMacro{\crpspCorSecundaria}
\DescribeMacro{\crpspCorAcessoria}
Cada publicação define sua paleta no preâmbulo com códigos HTML
(sem |#|):
\begin{verbatim}
\crpspCorPrincipal{004A8F}
\crpspCorSecundaria{E8452A}
\crpspCorAcessoria{D4A843}
\end{verbatim}
Os valores padrão são neutros (preto/cinza). As cores ficam disponíveis
no corpo do documento pelos aliases \texttt{cor1}, \texttt{cor2} e
\texttt{cor3} (para |\textcolor|, TikZ etc.).

\section{\cls{crpsp\_acessivel} — linha A5 (v0.2)}

Classe da linha editorial em A5 retrato, baseada em \cls{book}.
Comportamento documentado no \texttt{README.md} da pasta
\texttt{latex\_acessivel}; a v0.2 é um \emph{refactoring} sem mudança de
API: os blocos de engine, cores e fontspec migraram para
\pkg{crpsp-base}.

\DescribeMacro{[guia]}
\DescribeMacro{[livro]}
Exige exatamente uma opção de linha editorial:
|\documentclass[guia]{crpsp_acessivel}|. A linha \opt{livro} ainda não
está implementada (emite warning).

\section{\cls{crpsp-guia\_visual} — linha 16:9 (v0.1)}

Classe para guias visuais e apresentações institucionais em PDF
paisagem 16:9 ($338{,}67 \times 190{,}5$\,mm, padrão widescreen).
Baseada em \cls{article} — o documento é uma sequência plana de seções
e \emph{frames}, sem frontmatter/mainmatter/backmatter. Carrega
\pkg{crpsp-base} e \pkg{guia-visual}.

\DescribeMacro{[toc]}
Por padrão |\section| \emph{não} alimenta o sumário (a navegação é
feita pela navbar). A opção de classe \opt{toc} restaura o
|\addcontentsline| com o título longo.

\subsection{Estrutura do documento}

\DescribeMacro{\section}
\begin{verbatim}
\section[Título curto]{Título longo completo}
\end{verbatim}
Sintaxe idêntica à do Beamer. A seção: (1) inicia página nova;
(2) registra o título curto e a página de início no stream |.nav|;
(3) \emph{não} produz título visual no corpo — a seção é comunicada
pela barra de navegação. Sem o argumento opcional, o título longo é
usado também na navbar.

\DescribeEnv{frame}
Cada frame é uma unidade de layout autônoma que ocupa uma página
(termina com |\clearpage|). Opções em key--value:

\begin{center}
\begin{tabular}{ll}
\opt{layout=colunas} & padrão — o conteúdo usará o grid de 3 colunas \\
\opt{layout=full}    & largura total, sem grid \\
\opt{bg=<arquivo>}   & background de página inteira (via \pkg{eso-pic}) \\
\end{tabular}
\end{center}

\DescribeEnv{colunas}
\DescribeMacro{\coluna}
Grid de \textbf{3 colunas iguais}. Dentro de |colunas|, cada
|\coluna[span=N]{...}| ocupa $N$ unidades ($N \in \{1,2,3\}$); a soma
dos spans deve ser 3, senão o pacote emite warning. A ordem das
chamadas determina a ordem de leitura na árvore de tags
(seção~\ref{sec:ua2}).

\begin{verbatim}
\begin{frame}
  \begin{colunas}
    \coluna[span=2]{Conteúdo principal}%
    \coluna[span=1, c]{\includegraphics[width=\linewidth,
                      alt={Descrição da imagem}]{imagem.pdf}}%
  \end{colunas}
\end{frame}
\end{verbatim}

Termine cada argumento de |\coluna| com |%| após a chave para não vazar
espaços para o grid.

\paragraph{Alinhamento dentro da coluna.}
A coluna é uma caixa de \textbf{altura fixa} — por padrão a mancha do
frame. É isso que faz |\vfill| / |\vspace*{\fill}| funcionarem dentro
dela (numa caixa de altura natural a cola infinita não tem folga para
esticar e colapsa em zero) e que dá sentido às chaves de posição:

\begin{center}
\begin{tabular}{ll}
\opt{t} & \emph{(padrão)} conteúdo no topo \\
\opt{c} & centrado verticalmente \\
\opt{b} & rente ao rodapé da mancha \\
\opt{natural}      & altura natural (|\vfill| volta a ser inerte) \\
\opt{altura=<dim>} & altura explícita \\
\opt{justificado}  & \emph{(padrão)} texto justificado \\
\opt{esquerda}     & bloco alinhado à esquerda \\
\opt{centro}       & bloco centrado \\
\opt{direita}      & bloco alinhado à direita \\
\end{tabular}
\end{center}

As formas longas \opt{valign=t} (ou \texttt{c}, \texttt{b}) e
\opt{halign=centro} também são aceitas:

\begin{verbatim}
\coluna[span=1, b]{...}          % encosta no rodapé
\coluna[span=2, c, centro]{...}  % centrado nos dois eixos
\coluna[span=1]{Texto\vfill\includegraphics{fig}}  % figura ao pé
\end{verbatim}

Como a caixa tem altura fixa, conteúdo maior que a mancha passa a
gerar \emph{Overfull \cs{vbox}} — sinal legítimo de excesso de texto na
coluna; a saída é \opt{altura} ou \opt{natural}.

\subsection{Barra de navegação}\label{sec:nav}

Faixa em \texttt{cor1} nos primeiros 12\,mm do topo de cada página:
itens de seção da esquerda para a direita (ativo em fundo branco e
texto \texttt{cor1} negrito; inativos em branco), separadores verticais
e número de página à direita — referência visual: Relatório Anual 2025
do BNDES.

A seção ativa é \emph{inferida} por página: é a última seção cuja
página de início é $\le$ a página corrente. Os dados vêm de um
\textbf{stream \texttt{.nav} dedicado} (mesmo padrão do Beamer),
separado do |.aux| para não interagir com \pkg{hyperref}/toc/refs:
a primeira passagem grava um |\guia@navitem{índice}{título}{página}|
por seção; a segunda lê o arquivo em |\AtBeginDocument| (antes do
|\openout|, que trunca o arquivo).

Se a soma das larguras dos títulos (medidas em negrito, o estado mais
largo) exceder |\paperwidth|, o pacote avisa:
use |\section[Curto]{Longo}| para abreviar.

\DescribeMacro{\estilocorpo}
\DescribeMacro{\estiloaberturapagina}
\DescribeMacro{\estilovaziopagina}
Aliases públicos dos estilos de página \pkg{fancyhdr}
(\texttt{guia@corpo}, \texttt{guia@abertura}, \texttt{guia@vazio});
o último remove a navbar, para frames com background \emph{full-bleed}.

\subsection{Tipografia e listas}

Cadeia de fallback (a mesma da linha A5): Atkinson Hyperlegible Next
$\to$ TeX Gyre Heros $\to$ Computer Modern Sans. Parágrafos sem
indentação, separados por \textbf{uma entrelinha cheia}
(|\parskip| $=$ |\baselineskip|, rígido) — não é preciso inserir
|\vspace| entre parágrafos. O valor é fixado no hook
\texttt{begindocument/end} (antes disso |\baselineskip| ainda é 0\,pt)
e reaplicado dentro de cada |\coluna|, porque |\parbox| zera
|\parskip| via |\@parboxrestore|. Bullets e numeração de
listas em \texttt{cor1}, via redefinição de |\labelitemi| /
|\labelitemii| / |\labelenumi| (\emph{não} \pkg{enumitem} — ver
seção~\ref{sec:workarounds}).

\section{Conformidade PDF/UA-2}\label{sec:ua2}

O grid de colunas gera estrutura explícita: o ambiente |colunas| abre
um contêiner \texttt{Sect} e cada |\coluna| abre uma \texttt{Sect}
filha, na ordem do código-fonte — garantindo a ordem de leitura
coluna~1 $\to$ 2 $\to$ 3. \textbf{Não} se usa \pkg{multicol} nem
\pkg{paracol} (nenhum suporta tagged PDF).

Toda imagem informativa deve declarar texto alternativo:
\begin{verbatim}
\includegraphics[width=\linewidth,
  alt={Gráfico de barras mostrando...}]{figura.pdf}
\end{verbatim}

A navbar e o número de página são artefatos de paginação: o tagging é
suspenso (|\SuspendTagging|) durante sua construção.

\section{Workarounds phase-III}\label{sec:workarounds}

Registro dos contornos ativos nesta linha, descobertos com TeX Live
2026 (\pkg{tagpdf} 1.0c, \pkg{latex-lab} 2026-05). Antes de removê-los,
verificar o estado upstream com \texttt{monitor\_ctan.py}.

\begin{enumerate}
\item \textbf{\pkg{enumitem} é incompatível} com o módulo \pkg{block}
  do \pkg{latex-lab} (erro \emph{``Some keys specified on the itemize
  environment are unknown''}) — mesma família do problema conhecido do
  \pkg{titlesec}. Listas são estilizadas pelo mecanismo padrão do
  kernel.
\item \textbf{\pkg{tcolorbox} não fica inline} sob phase-III: o patch
  de tagging força |\par| em volta da box mesmo com
  \opt{nobeforeafter}, empilhando colunas. O grid usa |\parbox[t]|
  \emph{de altura fixa} — o mesmo workaround
  ``\,\emph{parbox duplo}\,'' da linha A5 para tabelas de layout.
\item \textbf{\cs{topskip} zerado no ambiente \texttt{colunas}}: no
  topo da página o \TeX{} insere cola de |\topskip| menos a altura da
  primeira caixa. Com a coluna de altura fixa essa altura é a da
  primeira linha ($\approx 5$\,pt), e os $\approx 5$\,pt de cola
  resultantes estouravam a mancha — uma página em branco extra por
  frame.
\item \textbf{\cs{settowidth} devolve 0\,pt dentro de
  \texttt{tikzpicture}} (interferência do módulo
  \pkg{latex-lab-testphase-tikz}). A navbar pré-mede todos os itens
  \emph{fora} do picture e só desenha dentro dele.
\item \textbf{\cs{SuspendTagging} em volta do \texttt{tikzpicture} do
  cabeçalho}: sem isso o módulo de tagging de TikZ tenta taggear a
  navbar como \texttt{Figure}, descarta o picture e deixa ``dados não
  processados'' (página extra \emph{Temporary page!}).
\item \textbf{Estruturas \texttt{Sect} não podem abrir dentro de
  parágrafo taggeado} (relação \texttt{P}$\to$\texttt{Sect} proibida
  no PDF 2.0): o grid desliga o \emph{paratagging} na linha de boxes
  (|\tagpdfparaOff|) e religa dentro de cada coluna.
\end{enumerate}

Dois contornos de implementação TikZ/pgf independentes de tagging:
não usar |\loop|/|\repeat| com |\node| dentro de \texttt{tikzpicture}
(o |\repeat| desbalanceia a pilha de condicionais do parser de paths;
usar |\foreach|), e não usar |\csuse| do \pkg{etoolbox} em expressões
de coordenada (o |\fi| do |\ifcsname| fica órfão no pgfmath; usar
|\csname| pré-expandido com |\edef|).

\section{Arquivos}

\begin{center}
\begin{tabular}{lll}
\textbf{arquivo} & \textbf{papel} & \textbf{versão} \\\hline
\pkg{crpsp-base.sty}          & infraestrutura compartilhada   & v0.1 \\
\cls{crpsp\_acessivel.cls}    & classe A5 (book)               & v0.2 \\
\pkg{guia-crpsp\_acessivel.sty} & layout da linha guia A5      & v0.1 \\
\cls{crpsp-guia\_visual.cls}  & classe 16:9 (article)          & v0.1 \\
\pkg{guia-visual.sty}         & layout da linha guia visual    & v0.1 \\
\texttt{exemplo-guia-visual.tex} & documento de teste 16:9     & --- \\
\texttt{guia-exemplo.tex}     & documento de teste A5          & --- \\
\end{tabular}
\end{center}

\end{document}
